% ---------------------------------------------------------------
% Trade-Offs
% Competing forces and the decisions made. Summarise / reference the
% ADRs that drove each significant trade-off.
% ---------------------------------------------------------------
\section{Trade-offs and Architectural Decisions}
\label{sec:tradeoffs}

In designing the TicketTailor platform, several competing architectural forces were balanced. Key decisions involved trade-offs between system performance, state consistency, operational complexity, and data privacy to satisfy the project's architecturally significant requirements (ASRs).

\subsection{Consistency vs. Low-Latency Scalability (RSVP Counts)}
To handle concurrent RSVP traffic spikes without saturating the primary relational database, read-path consistency was traded for low-latency scalability. Managing seat availability and attendee counts directly in PostgreSQL through transactional row-level updates or active aggregates (\texttt{COUNT(*)}) introduces severe lock contention and transaction failures under heavy write bursts. To prevent this, the live counter path is decoupled from PostgreSQL by utilizing a fast in-memory Redis cache line (\texttt{rsvp:count:\{event\_id\}}). Read operations for map browsing and event lists query the Redis counter directly.

This design introduces the risk of cache-to-database read drift. The Redis counter can temporarily diverge from the authoritative database rows if Redis experiences a crash or key eviction. However, the counter is only updated after a successful PostgreSQL commit, so a failed transaction never leaves a phantom counter. This drift was accepted because user-facing read updates do not require immediate transaction-level consistency. To safeguard correctness, the RSVP service executes a strict order of operations: the SQL row is inserted or deleted and durably committed to PostgreSQL before the write-through increment or decrement is sent to Redis. Seeding a cold key also runs before database writes to prevent double-counting or under-counting base values. Furthermore, a periodic background reconciliation task recalculates persisted RSVP rows in PostgreSQL and overwrites drifted Redis values. If Redis becomes entirely unreachable, the client connection wrapper raises connection timeouts quickly, allowing the API to degrade gracefully by falling back to SQL database queries instead of blocking the main ASGI event loop.

\subsection{Stateless Workers vs. Database Connection Pools}
To scale the notification engine under heavy bursts, the team faced a trade-off between database resource safety and notification boundary complexity. In a naive event-driven design, the notification worker (AWS Lambda) would receive thin event payloads and query the database to resolve recipient contact details and push subscriptions. However, fanning out updates concurrently to hundreds or thousands of club followers would spin up many parallel Lambda instances. This concurrent spike would quickly exhaust the PostgreSQL primary database connection pool, starving the core API and crashing the database.

To prevent connection starvation, the notification worker was designed to be completely database-free and stateless. Recipient resolution and message hydration are shifted upstream to the outbox relay running on ECS Fargate. The relay queries the database through module service interfaces to obtain emails, display names, and VAPID push keys, fanning out one fully pre-hydrated message per recipient to SNS, which fans out to the subscribed SQS queues. The trade-off is that fanning out hydrated messages widens the relay's dependencies and places personally identifiable information (PII) on the message bus. This security risk is mitigated by enabling server-side encryption with AWS managed keys (SSE-KMS) on the SNS topics and SQS queues. The database-free worker can scale rapidly without VPC or RDS connection bottlenecks, but because it has no database access, it cannot clean up expired push keys directly when push services return subscription errors. Instead, it logs these occurrences, leaving active database pruning as a potential improvement to be managed via secure API callbacks in a future version.

\subsection{At-Least-Once Delivery vs. Message Duplication}
To guarantee that no critical notification is permanently lost, reliability was prioritized over avoiding duplicate message deliveries. In a highly concurrent environment, ensuring exactly-once delivery across distributed processes (monolith API, outbox relay, SQS, and Lambda worker) is functionally impractical without heavy coordination overhead that degrades performance. The architecture uses a transactional outbox pattern combined with SQS retries to enforce an at-least-once delivery boundary.

The trade-off is that both the outbox relay and the SQS queues can emit duplicate messages under ordinary failure conditions. The relay uses an all-or-nothing publish strategy: all messages for a single outbox row must be published successfully before the row is marked as published, so a partial failure leaves the entire row unpublished for full retry on the next tick. Similarly, if the Lambda worker encounters a transient error (such as an email provider timeout) during a concurrent delivery run, it raises an exception and reports failures to SQS. SQS then redelivers the message, which may result in duplicate emails or push notifications on already-successful delivery channels. This duplicate overhead was accepted as a necessary compromise to ensure zero messages are lost. To mitigate the impact on users, the outbox relay generates deterministic message IDs based on the outbox row ID and recipient details, allowing downstream clients or worker handlers to identify and deduplicate incoming messages.

\subsection{Strict Module Boundaries vs. Development Friction}
To support future extensibility, the ease of ad-hoc database queries was traded for logical isolation and modular independence. In a traditional two-tier monolith, developers routinely write SQL queries that join tables across arbitrary modules or import classes globally, creating tightly coupled spaghetti code. To prevent this, strict schema-per-module separation is enforced, and physical database foreign keys across module schemas are prohibited. Module boundaries are also checked during CI runs using Import Linter.

This strict logical isolation introduces significant development friction. Common tasks, such as verifying club membership before allowing an event RSVP or retrieving event titles for organizer notifications, cannot use direct database joins or repository imports. Instead, modules must coordinate through narrow, service-level interfaces, registered shared callbacks, or multi-query aggregates. This structure increases the size of the codebase and the complexity of simple operations. However, the trade-off is justified because it prevents architectural degradation. The enforced boundaries ensure that the system maintains high cohesion and loose coupling, providing a clean extraction path where any single module can be migrated to a standalone microservice in a future version without refactoring the remaining codebase.